\documentclass[aspectratio=169,14pt]{beamer}
\usetheme{default}
\usepackage{fontspec}
\setmainfont{texgyrepagella-regular.otf}[BoldFont=texgyrepagella-bold.otf,ItalicFont=texgyrepagella-italic.otf,BoldItalicFont=texgyrepagella-bolditalic.otf]
\setsansfont{texgyreheros-regular.otf}[BoldFont=texgyreheros-bold.otf,ItalicFont=texgyreheros-italic.otf,BoldItalicFont=texgyreheros-bold.otf,BoldItalicFeatures={FakeSlant=0.2}]
\setmonofont{lmmono10-regular.otf}[Scale=0.85]
\setbeamertemplate{navigation symbols}{}
\usepackage{pgfpages}
\newif\ifpresenter
\ifdefined\PresenterVersion\presentertrue\else\presenterfalse\fi
\ifpresenter
  \setbeameroption{show notes on second screen=right}
  \setbeamertemplate{note page}[plain]
\fi
\usepackage{amsmath,amssymb,bm,booktabs,array,graphicx,xcolor}
\usepackage{tikz,pgfplots,adjustbox,listings,xurl}
\pgfplotsset{compat=1.18}
\usetikzlibrary{arrows.meta,positioning,calc,fit,shapes.geometric,decorations.pathreplacing,patterns}
% UIUC official brand colors, verified 2026-10-23.
% https://brand.illinois.edu/visual-identity/color
\definecolor{IlliniOrange}{HTML}{FF5F05}
\definecolor{IlliniBlue}{HTML}{13294B}
\definecolor{Storm}{HTML}{707372}
\definecolor{LightStorm}{HTML}{C6C7C6}
\definecolor{Industrial}{HTML}{1D58A7}
\definecolor{Patina}{HTML}{007E8E}
\tikzset{>=Stealth, every node/.style={font=\small,text=IlliniBlue}, box/.style={draw=IlliniBlue,rounded corners=2pt,align=center,inner sep=6pt,execute at begin node=\hyphenpenalty10000\relax}, orangebox/.style={box,draw=IlliniOrange,very thick}, edge/.style={->,thick,IlliniBlue}, accent/.style={->,very thick,IlliniOrange}}
\pgfplotsset{every axis/.append style={width=13.5cm,height=6.5cm,axis line style={IlliniBlue},tick label style={font=\small,text=IlliniBlue},label style={font=\small,text=IlliniBlue},legend style={font=\small,draw=none,fill=white},grid style={LightStorm},major grid style={LightStorm},axis lines=left,clip=false}, every axis plot/.append style={thick}}

\definecolor{slidepaper}{RGB}{234,220,201}
\setbeamercolor{background canvas}{bg=slidepaper}
\setbeamercolor{normal text}{fg=IlliniBlue}
\setbeamercolor{frametitle}{fg=IlliniBlue}
\setbeamercolor{title}{fg=IlliniBlue}
\setbeamercolor{item}{fg=IlliniOrange}
\setbeamerfont{frametitle}{size=\large,series=\bfseries}
\setbeamerfont{title}{size=\Large,series=\bfseries}
\setbeamertemplate{itemize item}{\color{IlliniOrange}$\bullet$}
\setbeamertemplate{itemize subitem}{\color{IlliniOrange}\footnotesize$\circ$}
\setbeamercolor{orangerule}{bg=IlliniOrange}
\setbeamertemplate{headline}{\begin{beamercolorbox}[wd=\paperwidth,ht=0.1cm,dp=0.1cm]{orangerule}\end{beamercolorbox}}
\setbeamertemplate{footline}{\begin{beamercolorbox}[wd=\paperwidth,ht=0.1cm,dp=0.1cm]{orangerule}\end{beamercolorbox}}
\setbeamersize{text margin left=0.65cm,text margin right=0.65cm}
\setlength{\leftmargini}{1.1em}
\newenvironment{slidebullets}{\begin{itemize}\setlength\itemsep{0.8em}}{\end{itemize}}
\pgfplotsset{every axis/.append style={legend style={fill=slidepaper}}}
\lstdefinelanguage{CoursePython}{keywords={from,import,for,in,if,else,def,return,assert,True,False,None},sensitive=true,morecomment=[l]{\#},morestring=[b]',morestring=[b]"}
\lstset{language=CoursePython,basicstyle=\ttfamily\small,keywordstyle=\color{IlliniBlue}\bfseries,commentstyle=\color{Storm},stringstyle=\color{IlliniBlue},showstringspaces=false,breaklines=true,columns=fullflexible,keepspaces=true,frame=none,aboveskip=0.7em,belowskip=0.7em}
\newcommand{\E}{\mathbb E}
\newcommand{\Var}{\operatorname{Var}}
\newcommand{\Cov}{\operatorname{Cov}}
\hypersetup{colorlinks=true,urlcolor=IlliniBlue,linkcolor=IlliniBlue}

\title{Marketing Analytics}
\subtitle{Integrated Decisions and Final Presentations}
\author{By Nathan Yang}
\date{Session 15}
\begin{document}
\begin{frame}
\titlepage
\end{frame}
\begin{frame}{An Integrated Marketing Decision}
\usebeamercolor[fg]{normal text}
\vfill
\begin{slidebullets}
\item Choose a feasible action and a defined customer population.
\item Connect measurement, modeling, and evaluation to that action.
\item State the costs, customer outcomes, and uncertainty that affect the choice.
\end{slidebullets}
\vfill
\note{The evidence chain can include a descriptive assessment, a predictive model, a causal comparison, and a scenario. Each component should retain its own scope and assumptions.}
\end{frame}
\begin{frame}{The Evidence Behind the Recommendation}
\usebeamercolor[fg]{normal text}
\vfill
\begin{slidebullets}
\item State the sample, observation unit, and outcome horizon.
\item Show the comparison that supports the main claim.
\item Explain which result changes the recommended action.
\end{slidebullets}
\vfill
\note{A result becomes decision-relevant when it changes a feasible choice under explicit costs or constraints. The audience should be able to assess the supporting comparison directly.}
\end{frame}
\begin{frame}{A Calculation the Audience Can Check}
\usebeamercolor[fg]{normal text}
\vfill
\begin{slidebullets}
\item Trace one figure to an aggregate table and source records.
\item Reproduce one prediction, coefficient interpretation, or policy value.
\item Retain complete code and data definitions with the result.
\end{slidebullets}
\vfill
\note{A hand-checkable calculation complements complete reproducible code. It provides independent evidence about the unit, denominator, and direction of the computation.}
\end{frame}
\begin{frame}{A Result That Could Change the Decision}
\usebeamercolor[fg]{normal text}
\vfill
\begin{slidebullets}
\item Identify the most consequential uncertain parameter.
\item Show the threshold at which the preferred action changes.
\item Consider a feasible test that could resolve that uncertainty.
\end{slidebullets}
\vfill
\note{Sensitivity should connect uncertainty to action. A wide parameter range is useful when it reveals whether the decision is stable or fragile.}
\end{frame}
\begin{frame}{Customer Consequences}
\usebeamercolor[fg]{normal text}
\vfill
\begin{slidebullets}
\item Report effects on access, time, value, and unwanted charges where relevant.
\item Distinguish voluntary continued use from difficulty leaving.
\item Explain how customers can refuse or correct the proposed interaction.
\end{slidebullets}
\vfill
\note{The relevant customer outcome depends on the product and action. A higher platform metric can conceal costs outside the optimized objective.}
\end{frame}
\begin{frame}{Regulatory Scope}
\usebeamercolor[fg]{normal text}
\vfill
\begin{slidebullets}
\item Identify the jurisdiction and covered conduct.
\item Separate the current rule from a proposal or allegation.
\item Match the evidence supporting a claim to what customers will understand.
\end{slidebullets}
\vfill
\note{The primary-source record should identify the effective date and consequential exclusions. A legal conclusion for one sector or jurisdiction should not be transferred mechanically to another.}
\end{frame}
\begin{frame}{Presentation Discussion}
\usebeamercolor[fg]{normal text}
\vfill
\begin{slidebullets}
\item Is the recommendation supported for the stated population?
\item What analytical or customer risk has the largest consequence?
\item Which part of the result can be reproduced independently?
\end{slidebullets}
\vfill
\note{Feedback should name an observable gap in the evidence or a specific counterexample. The strongest response supplies a corrected comparison or explains why the original remains valid.}
\end{frame}
\begin{frame}{A Deployment Comparison}
\usebeamercolor[fg]{normal text}
\vfill
\begin{slidebullets}
\item Retain a benchmark and a defined evaluation population.
\item Monitor outcomes and data coverage after the action changes behavior.
\item Set a response when quality or customer outcomes deteriorate.
\end{slidebullets}
\vfill
\note{A policy that changes exposure can also change future training data. Evaluation after deployment needs versioned inputs, outcomes, and policy rules.}
\end{frame}
\begin{frame}{An Evidence-Based Recommendation}
\usebeamercolor[fg]{normal text}
\vfill
\begin{slidebullets}
\item Name the action, expected consequence, and decision threshold.
\item Keep consequential uncertainty beside the estimate.
\item Identify the next concrete extension supported by the evidence.
\end{slidebullets}
\vfill
\note{A conditional recommendation can be decisive when its conditions are measurable. The final decision should remain traceable to the data, assumptions, and customer outcomes.}
\end{frame}
\end{document}
